%! Unit 5: Rational Functions — Summative Assessment — Practice Test --- Study Copy (KEY)
\documentclass[10pt]{article}
\usepackage{algebra2-article}
\usepackage{algebra2-key}   % pulls in -boxes; do NOT also load -boxes
% Coordinate grid: {xmin}{xmax}{ymin}{ymax}
\newcommand{\lgrid}[4]{%
  \draw[step=1, linegray!40, very thin] (#1,#3) grid (#2,#4);
  \draw[->, semithick, charcoal] (#1,0)--(#2,0) node[below right, font=\scriptsize]{$x$};
  \draw[->, semithick, charcoal] (0,#3)--(0,#4) node[left, font=\scriptsize]{$y$};}
\newcommand{\dpt}[2]{\fill[charcoal] (#1,#2) circle (3pt);}

% Part-divider strip (headlinebox takes one arg: the background color)
\newcommand{\parthead}[1]{%
  \vspace{6pt}\noindent
  \begin{headlinebox}{forest}\color{white}\bfseries #1\end{headlinebox}%
  \vspace{4pt}}

\begin{document}

\pageheader{Unit 5: Rational Functions}{Practice Test --- Study Copy --- Answer Key}
\namedateperiod

\vspace{0.10in}
\begin{remindbox}
\small
\textbf{This is a practice test.} It mirrors the real Unit 5 test in length, format,
and the ideas it covers, but the numbers and contexts differ. Work every problem
with no notes, then check your answers against the key.
\end{remindbox}

% ======================================================================
\parthead{Part A --- Vocabulary (8 pts)}
Write the letter of the best description next to each term.

\vspace{4pt}
\noindent
\begin{minipage}[t]{0.40\textwidth}
\begin{enumerate}[leftmargin=2.2em, itemsep=7pt, label=\arabic*.]
  \item Rational function \ans{D}
  \item Domain restriction \ans{A}
  \item Vertical asymptote \ans{H}
  \item Hole \ans{E}
  \item Horizontal asymptote \ans{B}
  \item Least common denominator \ans{G}
  \item Extraneous solution \ans{F}
  \item Constant of variation \ans{C}
\end{enumerate}
\end{minipage}%
\hfill
\begin{minipage}[t]{0.57\textwidth}
\small
\begin{description}[leftmargin=1.4em, itemsep=3pt]
  \item[(A)] A value that makes the \emph{original} denominator zero, so it can never be an input.
  \item[(B)] A horizontal line the graph approaches at its far left and far right; found by comparing the degrees of the numerator and denominator.
  \item[(C)] The number $k$ in $y=kx$ or $y=\frac{k}{x}$; you find it from one known pair of values.
  \item[(D)] A quotient of two polynomials whose denominator is not zero.
  \item[(E)] A single missing point on an otherwise unbroken curve, left behind by a factor that cancels.
  \item[(F)] A number that turns up when you solve but must be thrown out, because it is an excluded value of the original equation.
  \item[(G)] The smallest expression that every denominator divides evenly, built from the factored denominators.
  \item[(H)] A vertical line the graph runs alongside but never touches, left by a factor that does \emph{not} cancel.
\end{description}
\end{minipage}

% ======================================================================
\parthead{Part B --- Multiple Choice (12 pts)}
Circle the best answer.

\begin{enumerate}[leftmargin=1.9em, itemsep=9pt]
  \item The domain of $f(x)=\dfrac{x+5}{x^2-9}$ is all real numbers except
    \hfill (A) $x=-5$ \quad (B) $x=3$ \quad \textbf{(C) $x=\pm 3$} \ans{$\leftarrow$} \quad (D) none --- all reals
  \item Simplified completely, $\dfrac{x^2-25}{x^2+3x-10}$ is
    \hfill (A) $\dfrac{x-5}{x-2}$, $x\ne 2$ \quad \textbf{(B) $\dfrac{x-5}{x-2}$, $x\ne -5,\,2$} \ans{$\leftarrow$} \quad (C) $\dfrac{x+5}{x-2}$, $x\ne 2$ \quad (D) $\dfrac{-5}{3x-10}$
  \item The horizontal asymptote of $f(x)=\dfrac{3x^2-1}{x^2+4x}$ is
    \hfill (A) $y=0$ \quad (B) $y=\frac13$ \quad \textbf{(C) $y=3$} \ans{$\leftarrow$} \quad (D) there is none
  \item Which statement about $h(x)=\dfrac{x^2-x-6}{x^2-9}$ is true?
    \\[3pt]\hspace*{1.2em} (A) There are vertical asymptotes at $x=3$ and at $x=-3$.
    \\[2pt]\hspace*{1.2em} \textbf{(B) There is a hole at $x=3$ and a vertical asymptote at $x=-3$.} \ans{$\leftarrow$}
    \\[2pt]\hspace*{1.2em} (C) There is a hole at $x=-3$ and a vertical asymptote at $x=3$.
    \\[2pt]\hspace*{1.2em} (D) There is a hole at $x=-2$ and a vertical asymptote at $x=3$.
  \item The graph of $g(x)=\dfrac{1}{x-4}+2$ has asymptotes
    \hfill \textbf{(A) $x=4$, $y=2$} \ans{$\leftarrow$} \quad (B) $x=-4$, $y=2$ \quad (C) $x=4$, $y=0$ \quad (D) $x=2$, $y=4$
  \item $y$ varies \emph{inversely} as $x$, and $y=8$ when $x=3$. When $x=12$, $y=$
    \hfill \textbf{(A) $2$} \ans{$\leftarrow$} \quad (B) $6$ \quad (C) $24$ \quad (D) $32$
\end{enumerate}

\begin{teachernote}
Item 1: set the \emph{denominator} to zero --- $x^2-9=0$ gives $x=\pm3$; (A) is the numerator-zero
error and (B) keeps only one of the two. Item 2: $\frac{(x-5)(x+5)}{(x+5)(x-2)}$. (A) is the
signature Unit 5 error --- restrictions read off the \emph{simplified} denominator, losing the hole
at $x=-5$; (D) cancels the $x^2$ terms. Item 3: equal degrees, so the horizontal asymptote is the
ratio of leading coefficients, $\frac31$; (A) is the $n<m$ rule, (B) inverts the ratio. Item 4:
$\frac{(x-3)(x+2)}{(x-3)(x+3)}$ --- the $(x-3)$ cancels (hole), the $(x+3)$ does not (wall). (A)
skips the cancel test entirely; (D) reads a restriction off the numerator. Item 5: $x-4$ shifts
\emph{right}, so the wall moves to $x=4$; $+2$ lifts the horizontal asymptote to $y=2$. Item 6:
$k=xy=24$, so $y=\frac{24}{12}=2$. (D) is the direct-variation answer, (C) reports $k$ itself.
\end{teachernote}

% ======================================================================
\parthead{Part C --- Short Answer \& Computation (40 pts)}
Show your work. State every restriction on the variable.

\begin{enumerate}[leftmargin=1.9em, itemsep=10pt]
  \item \textbf{Read the graph.} The rational function $f$ is graphed below; the dashed lines are its
        asymptotes and each gridline is one unit.
        \begin{center}
        \begin{tikzpicture}[scale=0.45]
          \lgrid{-6}{6}{-6}{6}
          \draw[navy, dashed, semithick] (2,-6)--(2,6);
          \draw[navy, dashed, semithick] (-6,1)--(6,1);
          \begin{scope}
            \clip (-6,-6) rectangle (6,6);
            \draw[very thick, forest, domain=-6:1.14, samples=160, smooth]
              plot (\x,{((\x)+4)/((\x)-2)});
            \draw[very thick, forest, domain=3.2:6, samples=160, smooth]
              plot (\x,{((\x)+4)/((\x)-2)});
          \end{scope}
          \foreach \x in {-4,-2,2,4}
            \draw[charcoal] (\x,-0.15)--(\x,0.15) node[below=1pt, font=\tiny]{$\x$};
          \dpt{-4}{0}
          \dpt{0}{-2}
        \end{tikzpicture}
        \end{center}
        (a) Equation of the vertical asymptote: \ans{$x=2$} \quad
            Equation of the horizontal asymptote: \ans{$y=1$}
        \\[4pt]
        (b) $x$-intercept: \ans{$(-4,0)$}; \quad $y$-intercept: \ans{$(0,-2)$}
        \\[4pt]
        (c) Domain (interval notation): \ans{$(-\infty,2)\cup(2,\infty)$}
        \\[4pt]
        (d) Range (interval notation): \ans{$(-\infty,1)\cup(1,\infty)$}
        \\[4pt]
        (e) End behavior: as $x\to\pm\infty$, $f(x)\to$ \ans{$1$}
        \\[4pt]
        (f) Interval(s) where $f$ is decreasing: \ans{$(-\infty,2)$ and $(2,\infty)$ --- decreasing on \emph{each} branch, but not on their union}
        \\[4pt]
        (g) Interval(s) where $f(x)>0$: \ans{$(-\infty,-4)\cup(2,\infty)$}

  \item \textbf{Simplify and state the restrictions.} Read the restrictions from the \emph{original}
        denominator.
        \\[4pt]
        (a) $\dfrac{3x^2-27}{x^2+x-12}$ \hspace{2.5cm} (b) $\dfrac{x^2-16}{4-x}$
        \\[4pt]
        \ans{(a) GCF first, then difference of squares:
        $\dfrac{3(x-3)(x+3)}{(x+4)(x-3)}=\dfrac{3(x+3)}{x+4}$, \ $x\ne 3,\ x\ne -4$. The
        restriction $x\ne3$ survives even though the factor cancelled --- that is the hole.\\[3pt]
        (b) $4-x=-(x-4)$, so $\dfrac{(x-4)(x+4)}{-(x-4)}=-(x+4)$, \ $x\ne 4$. Opposite binomials
        leave a factor of $-1$ behind.}

  \item \textbf{Multiply or divide.} Factor first; give the simplest form with its restrictions.
        \\[6pt]
        (a) $\dfrac{x^2-4}{x^2+6x+9}\cdot\dfrac{x+3}{x-2}$
        \ans{$=\dfrac{(x-2)(x+2)}{(x+3)^2}\cdot\dfrac{x+3}{x-2}=\dfrac{x+2}{x+3}$, \ $x\ne -3,\ x\ne 2$.
        Cancelling runs across the whole product --- the $(x-2)$ in the first numerator meets the
        $(x-2)$ in the second denominator.}
        \\
        (b) $\dfrac{x^2-x-6}{x^2-1}\div\dfrac{x-3}{x+1}$
        \ans{Keep--change--flip: $\dfrac{(x-3)(x+2)}{(x-1)(x+1)}\cdot\dfrac{x+1}{x-3}
        =\dfrac{x+2}{x-1}$, \ $x\ne 1,\ x\ne -1,\ x\ne 3$. Three sources of restriction: the original
        denominator ($\pm1$), the divisor's denominator ($-1$ again), and the divisor's
        \emph{numerator} ($x=3$, because it becomes a denominator after the flip).}

  \item \textbf{Add, subtract, and simplify a complex fraction.}
        \\[6pt]
        (a) $\dfrac{5}{x+2}-\dfrac{3}{x-1}$
        \ans{LCD $=(x+2)(x-1)$:
        $\dfrac{5(x-1)-3(x+2)}{(x+2)(x-1)}=\dfrac{5x-5-3x-6}{(x+2)(x-1)}
        =\dfrac{2x-11}{(x+2)(x-1)}$, \ $x\ne -2,\ x\ne 1$. The subtraction hits \emph{both} terms of
        $3(x+2)$.}
        \\
        (b) $\dfrac{x}{x^2-9}+\dfrac{2}{x+3}$
        \ans{Factor first: $x^2-9=(x-3)(x+3)$, so the LCD is already $(x-3)(x+3)$:
        $\dfrac{x+2(x-3)}{(x-3)(x+3)}=\dfrac{3x-6}{(x-3)(x+3)}=\dfrac{3(x-2)}{(x-3)(x+3)}$,
        \ $x\ne \pm 3$. Nothing cancels --- $(x-2)$ is not a factor of either denominator factor.}
        \\
        (c) $\dfrac{\ \frac{1}{x}+\frac{1}{2}\ }{\ \frac{1}{x}-\frac{1}{2}\ }$
        \ans{Combine top and bottom over the common denominator $2x$, then divide:
        $\dfrac{\frac{2+x}{2x}}{\frac{2-x}{2x}}=\dfrac{2+x}{2x}\cdot\dfrac{2x}{2-x}=\dfrac{2+x}{2-x}$,
        \ $x\ne 0,\ x\ne 2$. The $x\ne0$ comes from the little fractions and would be invisible in
        the final answer.}

  \item \textbf{The parent function and its transformations.} Let $g(x)=\dfrac{1}{x+3}-2$.
        \\[4pt]
        (a) Describe the transformation(s) that take $y=\dfrac1x$ to $g$. \ans{left 3 units, down 2 units}
        \\[4pt]
        (b) Vertical asymptote: \ans{$x=-3$}; \quad horizontal asymptote: \ans{$y=-2$}
        \\[4pt]
        (c) Domain: \ans{$x\ne -3$}; \quad range: \ans{$y\ne -2$}
        \\[4pt]
        (d) A different hyperbola of the form $y=\dfrac{1}{x-h}+k$ has vertical asymptote $x=1$ and
            horizontal asymptote $y=4$. Write its equation. \ans{$y=\dfrac{1}{x-1}+4$}

  \item \textbf{Analyze from the equation.} For $f(x)=\dfrac{2x^2-8}{x^2-x-6}$:
        \\[4pt]
        (a) Factor the numerator and denominator, and list every excluded value.
        \ans{$f(x)=\dfrac{2(x-2)(x+2)}{(x-3)(x+2)}$. The original denominator is zero at $x=3$ and
        $x=-2$, so both are excluded.}
        \\
        (b) Which excluded value gives a \emph{hole} and which gives a \emph{vertical asymptote}?
            Give the coordinates of the hole.
        \ans{$(x+2)$ cancels, so $x=-2$ is a hole; $(x-3)$ stays, so $x=3$ is a vertical asymptote.
        Height of the hole from the simplified form $\dfrac{2(x-2)}{x-3}$ at $x=-2$:
        $\dfrac{2(-4)}{-5}=\dfrac85$, so the hole is $\left(-2,\tfrac85\right)$.}
        \\
        (c) Horizontal asymptote: \ans{$y=2$} \quad (Justify with a degree comparison.)
        \ans{Numerator and denominator are both degree 2, so the asymptote is the ratio of leading
        coefficients, $\tfrac21=2$.}
        \\[4pt]
        (d) $x$-intercept: \ans{$(2,0)$}; \quad $y$-intercept: \ans{$\left(0,\tfrac43\right)$}

  \item \textbf{Solve. Check every answer against the excluded values.}
        \\[6pt]
        (a) $\dfrac{x}{x+2}+\dfrac{1}{x-1}=1$
        \ans{Excluded: $x\ne -2,\,1$. Multiply by the LCD $(x+2)(x-1)$:
        $x(x-1)+(x+2)=(x+2)(x-1)$, so $x^2+2=x^2+x-2$, giving $x=4$. Since $4$ is not excluded,
        $x=4$ is the solution.}
        \\
        (b) $\dfrac{x}{x-4}=\dfrac{4}{x-4}+3$
        \ans{Excluded: $x\ne 4$. Multiply by $(x-4)$: $x=4+3(x-4)=3x-8$, so $-2x=-8$ and $x=4$. But
        $4$ is exactly the excluded value --- it is \textbf{extraneous}. The equation has
        \textbf{no solution}.}
        \\
        (c) $\dfrac{x}{x-2}+\dfrac{2}{x}=\dfrac{4}{x^2-2x}$
        \ans{Factor the last denominator: $x^2-2x=x(x-2)$, so the LCD is $x(x-2)$ and the excluded
        values are $x\ne 0,\,2$. Multiplying through: $x^2+2(x-2)=4$, so $x^2+2x-8=0$ and
        $(x+4)(x-2)=0$, giving $x=-4$ or $x=2$. Throw out $x=2$ (excluded). Only $x=-4$ survives;
        check: $\tfrac{-4}{-6}+\tfrac{2}{-4}=\tfrac23-\tfrac12=\tfrac16$ and
        $\tfrac{4}{24}=\tfrac16$ \checkmark}

  \item \textbf{Variation.} The table below records a pair of related quantities.
        \begin{center}
        \renewcommand{\arraystretch}{1.25}
        \begin{tabular}{|c|c|c|c|c|}
        \hline
        $x$ & $2$ & $4$ & $5$ & $10$ \\ \hline
        $y$ & $30$ & $15$ & $12$ & $6$ \\ \hline
        \end{tabular}
        \end{center}
        (a) Is $y$ directly or inversely proportional to $x$? Give the computation from the table that
            proves it.
        \\[2pt]
        \ans{Inversely --- the \emph{products} are constant: $2\cdot30=4\cdot15=5\cdot12=10\cdot6=60$ (the ratios $y/x$ are not)}
        \\[4pt]
        (b) Find the constant of variation $k$ and write the equation. \ans{$k=60$, so $y=\dfrac{60}{x}$}
        \\[4pt]
        (c) Use your equation to find $y$ when $x=8$. \ans{$y=\dfrac{60}{8}=7.5$}
\end{enumerate}

\begin{teachernote}
Item 1(f) is the Unit 5 habit-breaker: the answer is two \emph{separate} intervals, never
$(-\infty,\infty)$ minus a point written as one interval --- nothing crosses a wall. Item 1(d):
accept ``all reals except $y=1$.'' Items 2--4: a simplified form with \emph{no} restriction list is
half an answer --- score it as such; the restriction must come from the \emph{original} denominator.
Item 3(b) is the three-sources item: watch for the missing $x\ne3$ from the divisor's numerator.
Item 6(b): the hole's height must come from the \emph{simplified} form; substituting into the
original gives $\tfrac00$. Item 7(b) is the payoff loop --- ``$x=4$'' as a final answer earns no
credit; the check against the excluded values is the graded step. Accept ``no solution'' or
``$\varnothing$''. Item 8(a): a student who computes $y/x$ and finds it changing has still shown
the right work, provided they then check the products.
\end{teachernote}

% ======================================================================
\parthead{Part D --- Extended Response (12 pts)}
Explain your reasoning in full sentences.

\begin{enumerate}[leftmargin=1.9em, itemsep=16pt]
  \item \textbf{The full analysis.} Consider $f(x)=\dfrac{x^2-5x+4}{x^2-16}$.
        \textbf{(a)} Factor both parts and list every excluded value of the \emph{original} function.
        \textbf{(b)} Run the cancel test: which excluded value is a hole and which is a vertical
        asymptote? Give the hole's coordinates and explain how you found its height.
        \textbf{(c)} State the horizontal asymptote, justify it with a degree comparison, and write the
        end behavior it describes.
        \textbf{(d)} Give the $x$- and $y$-intercepts.
        \textbf{(e)} Write the domain in interval notation.
        \textbf{(f)} A classmate simplifies $f$ to $\dfrac{x-1}{x+4}$ and says the two functions are the
        same. Say exactly where they agree and where they do not.
        \ans{(a) $f(x)=\dfrac{(x-1)(x-4)}{(x-4)(x+4)}$. The original denominator is zero at $x=4$ and
        $x=-4$, so both are excluded.\\[3pt]
        (b) $(x-4)$ cancels, so $x=4$ is a \emph{hole}; $(x+4)$ does not, so $x=-4$ is a
        \emph{vertical asymptote}. The hole's height comes from the simplified form
        $\dfrac{x-1}{x+4}$ evaluated at $x=4$: $\dfrac{3}{8}$. The hole is
        $\left(4,\tfrac38\right)$. Substituting $4$ into the \emph{original} would give $\tfrac00$,
        which is why the simplified form is the one to use.\\[3pt]
        (c) Both degrees are $2$, so the horizontal asymptote is the ratio of leading coefficients,
        $\tfrac11=1$: the line $y=1$. That is the end-behavior statement --- as $x\to-\infty$ and as
        $x\to+\infty$, $f(x)\to 1$.\\[3pt]
        (d) $x$-intercept: the simplified numerator is zero at $x=1$, giving $(1,0)$. $y$-intercept:
        $f(0)=\dfrac{4}{-16}=-\dfrac14$, giving $\left(0,-\tfrac14\right)$.\\[3pt]
        (e) $(-\infty,-4)\cup(-4,4)\cup(4,\infty)$.\\[3pt]
        (f) They agree at every input both accept --- every $x$ other than $4$ and $-4$. They do not
        agree at $x=4$: the simplified expression happily returns $\tfrac38$ there, but $f(4)$ does
        not exist. Two equivalent expressions can have different \emph{domains}, so cancelling never
        restores an input.}

  \item \textbf{Modeling.} A club pays a one-time setup fee of \$180 to a screen printer, plus \$7 for
        each shirt printed. The \emph{average cost per shirt} for an order of $n$ shirts is
        \[ A(n)=\frac{7n+180}{n}. \]
        \textbf{(a)} Explain why $A$ is a rational function, and state the domain that makes sense for
        this situation.
        \textbf{(b)} Find $A(20)$ and $A(60)$, with units. What is happening to the average cost?
        \textbf{(c)} Rewrite $A(n)$ in the form $7+\dfrac{180}{n}$ and use it to state the horizontal
        asymptote. What does that number mean about this order?
        \textbf{(d)} Find the order size for which the average cost is exactly \$9 per shirt.
        \textbf{(e)} Can the average cost ever be exactly \$7? Answer using the horizontal asymptote.
        \ans{(a) It is a quotient of two polynomials, $7n+180$ over $n$. Algebraically the only
        restriction is $n\ne 0$; in context $n$ must be a positive whole number of shirts, so the
        sensible domain is $n=1,2,3,\dots$ --- an order of $0$ shirts has no ``per shirt'' cost at
        all.\\[3pt]
        (b) $A(20)=\dfrac{140+180}{20}=\dfrac{320}{20}=\$16$ per shirt;
        $A(60)=\dfrac{420+180}{60}=\dfrac{600}{60}=\$10$ per shirt. The average cost is
        \emph{falling} as the order grows, because the one-time \$180 is being split among more
        shirts.\\[3pt]
        (c) Divide term by term: $A(n)=\dfrac{7n}{n}+\dfrac{180}{n}=7+\dfrac{180}{n}$. As
        $n\to\infty$ the second piece shrinks toward $0$, so the horizontal asymptote is $y=7$. That
        \$7 is the true printing cost of one shirt --- the setup fee per shirt is the whole
        difference between the average cost and \$7.\\[3pt]
        (d) Solve $\dfrac{7n+180}{n}=9$. Multiply by $n$ (allowed, since $n\ne0$): $7n+180=9n$, so
        $180=2n$ and $n=90$ shirts.\\[3pt]
        (e) No. $A(n)=7+\dfrac{180}{n}$, and $\dfrac{180}{n}$ is positive for every positive $n$, so
        $A(n)>7$ always. The graph approaches the line $y=7$ but never reaches it --- the setup fee
        never fully disappears, it only gets spread thinner.}
\end{enumerate}

\begin{teachernote}
\textbf{Scoring Part D (6 pts each).} D1: 1 pt (a) both factorings with \emph{both} excluded values
from the original denominator; 2 pts (b) hole at $x=4$ \emph{and} wall at $x=-4$, with the height
$\tfrac38$ taken from the simplified form (1 pt only if the two are named but the height is missing
or computed from the original); 1 pt (c) $y=1$ with the degree comparison, not just the answer;
1 pt (d) both intercepts; 1 pt (e) the three-piece union --- a domain written as one interval, or
missing the $x\ne4$, loses it. (f) is the concept check: fold it into (b)/(e) credit if you prefer,
but a student who says the two functions are simply ``the same'' has missed the unit's central idea.
D2: 1 pt (a) rational \emph{and} a context domain of positive whole numbers; 1 pt (b) both values
with dollar units and the observation that cost is falling; 2 pts (c) the rewrite \emph{and}
$y=7$ interpreted as the real per-shirt printing cost (1 pt for $y=7$ alone); 1 pt (d) $n=90$;
1 pt (e) ``no,'' justified by $\tfrac{180}{n}>0$ or by the asymptote never being reached. Do not
accept ``because it's an asymptote'' with no reason attached.
\end{teachernote}

\end{document}
